% source: Robin Hartshorne, "Algebraic Geometry", GTM 52 (1977)
% pdf_page: 0147
% doc_page: 130 (Chapter II, Section 6, "Divisors")
% retrieved: 2026-07-26
% transcribed_by: codex/gpt-5 (vision, rendered at 200dpi via pdftoppm)

% running head: II Schemes / p. 130

\DeclareMathOperator{\Div}{Div}

% [continuation of the plane-curve example from the previous page]

projective space $(\mathbf{P}^2)^*$. We call this set of divisors a linear
system of divisors on $C$. Note that the embedding of $C$ in $\mathbf{P}^2$
can be recovered just from knowing this linear system: if $P$ is a point of
$C$, we consider the set of divisors in the linear system which contain $P$.
They correspond to the lines $L \in (\mathbf{P}^2)^*$ passing through $P$,
and this set of lines determines $P$ uniquely as a point of $\mathbf{P}^2$.
This connection between linear systems and embeddings in projective space
will be studied in detail in \S 7.

This example should already serve to illustrate the importance of divisors.
To see the relation among the different divisors in the linear system, let
$L$ and $L'$ be two lines in $\mathbf{P}^2$, and let $D = L \cap C$ and
$D' = L' \cap C$ be the corresponding divisors. If $L$ and $L'$ are defined
by linear homogeneous equations $f = 0$ and $f' = 0$ in $\mathbf{P}^2$, then
$f/f'$ gives a rational function on $\mathbf{P}^2$, which restricts to a
rational function $g$ on $C$. Now by construction, $g$ has zeros at the
points of $D$, and poles at the points of $D'$, counted with multiplicities,
in a sense which will be made precise below. We say that $D$ and $D'$ are
linearly equivalent, and the existence of such a rational function can be
taken as an intrinsic definition of the linear equivalence. We will make
these concepts more precise in our formal discussion, starting now.

\subsection*{Weil Divisors}

\begin{definition}
We say a scheme $X$ is \emph{regular in codimension one} (or sometimes
\emph{nonsingular in codimension one}) if every local ring
$\mathcal{O}_x$ of $X$ of dimension one is regular.
\end{definition}

The most important examples of such schemes are nonsingular varieties over a
field (I, \S 5) and noetherian normal schemes. On a nonsingular variety the
local ring of every closed point is regular (I, 5.1), hence all the local
rings are regular, since they are localizations of the local rings of closed
points. On a noetherian normal scheme, any local ring of dimension one is an
integrally closed domain, hence is regular (I, 6.2A).

In this section we will consider schemes satisfying the following condition:
\[
  \begin{aligned}
    (*)\quad &X \text{ is a noetherian integral separated scheme which is
      regular in} \\
    &\text{codimension one.}
  \end{aligned}
\]

\begin{definition}
Let $X$ satisfy $(*)$. A \emph{prime divisor} on $X$ is a closed integral
subscheme $Y$ of codimension one. A \emph{Weil divisor} is an element of the
free abelian group $\Div X$ generated by the prime divisors. We write a
divisor as $D = \sum n_iY_i$, where the $Y_i$ are prime divisors, the $n_i$
are integers, and only finitely many $n_i$ are different from zero. If all
the $n_i \ge 0$, we say that $D$ is \emph{effective}.
\end{definition}

If $Y$ is a prime divisor on $X$, let $\eta \in Y$ be its generic point.
Then the local ring $\mathcal{O}_{\eta,X}$ is a discrete valuation ring with
quotient field $K$, the function field of $X$. We call the corresponding
discrete valuation $v_Y$ the \emph{valuation of $Y$}. Note that since $X$ is
separated, $Y$ is uniquely deter-
% [continued on page-0148.tex]
